\section{Introduction}
\label{sec:introduction}

% Outline:
% 1. Why adaptive traffic signal control matters (congestion, delay, emissions).
% 2. Existing approaches: fixed-time and actuated plans, max-pressure, and RL controllers. RL needs long,
%    scenario-specific training and generalizes poorly to unseen intersections and demand.
% 3. Generative LLM controllers exist but are slow and nondeterministic. CLMs are a new class of model that ranks a
%    fixed set of candidate actions against a state instead of generating text, which suits signal control, where
%    the candidate set (green phases) is small and stable.
% 4. This paper: first study of CLMs for signal control, comparing the base and a traffic-specialized CLM-8B against
%    fixed-time, max-pressure, Q-learning and deep RL baselines, evaluated on held-out traffic.
% 5. Research questions RQ1--RQ5 (keep in sync with RESEARCH_TODO.md).
% 6. Contributions list.

\todo[inline]{Draft the introduction following the outline in the source comments.}

The research questions are:
\begin{itemize}
  \item \textbf{RQ1:} Can the \clmbase, used zero-shot, outperform fixed-time control and match max-pressure?
  \item \textbf{RQ2:} Does a \clmtraffic{} outperform the \clmbase, and does it match or outperform Q-learning and
        deep RL?
  \item \textbf{RQ3:} How much data does each method need to reach its performance?
  \item \textbf{RQ4:} Do the CLM controllers generalize to unseen intersections and demand patterns?
  \item \textbf{RQ5:} Is CLM inference fast enough for real-time signal control?
\end{itemize}

\todo[inline]{Contributions list.}
